%% kapa: exam=1 round=29 year=2018 period=2 form=A subject=accounting subject_ko=회계학 seq=12 number=52
%% source: Q-Net 감정평가사 기출문제 (cst00309) artlSeq=5202808 file=extracted/1008/5202808/2018년 제29회 감정평가사 1차시험 기출문제/2018년 제29회 감정평가사 1차시험 기출문제/2교시 A형.pdf
%% conversion: kordoc 4.24.51 -> tools/gen1.py (auto); review=proofread-sonnet; flags=none
\begin{question}{52}
다음은 (주)감평이 채택하고 있는 확정급여제도와 관련한 자료이다.
\begin{framedtext}
\begin{tabularx}{\linewidth}{@{}Xr@{}}
○ 확정급여채무 계산 시 적용하는 할인율 & 연 5\% \\
○ 기초 확정급여채무의 현재가치 & ￦700,000 \\
○ 기초 사외적립자산의 공정가치 & 600,000 \\
○ 당기근무원가 & 73,000 \\
○ 사외적립자산에 대한 기여금 출연(기말 납부) & 90,000 \\
○ 퇴직급여 지급액(사외적립자산에서 기말 지급) & 68,000 \\
○ 기말 사외적립자산의 공정가치 & 670,000 \\
○ 기말 재무상태표에 표시된 순확정급여부채 & 100,000 \\
\end{tabularx}
\end{framedtext}
(주)감평의 확정급여제도 적용이 포괄손익계산서의 당기순이익과 기타포괄이익에 미치는 영향은 각각 얼마인가? \par
\hspace{2.5em}\makebox[6cm][l]{\uline{당기순이익에 미치는 영향}}\uline{기타포괄이익에 미치는 영향} \par
\begin{choices}
  \item \makebox[6cm][l]{￦108,000 감소}￦48,000 감소
  \item \makebox[6cm][l]{￦108,000 감소}￦48,000 증가
  \item \makebox[6cm][l]{￦108,000 감소}￦12,000 감소
  \item \makebox[6cm][l]{￦78,000 감소}￦12,000 증가
  \item \makebox[6cm][l]{￦78,000 감소}￦12,000 감소
\end{choices}
\end{question}
\answer{5}
